
\documentclass{article}
\usepackage{colortbl}
\usepackage{makecell}
\usepackage{multirow}
\usepackage{supertabular}

\begin{document}

\newcounter{utterance}

\centering \large Interaction Transcript for game `codenames', experiment `risk\_low', episode 2 with claude{-}opus{-}4{-}8{-}azure.
\vspace{24pt}

{ \footnotesize  \setcounter{utterance}{1}
\setlength{\tabcolsep}{0pt}
\begin{supertabular}{c@{$\;$}|p{.15\linewidth}@{}p{.15\linewidth}p{.15\linewidth}p{.15\linewidth}p{.15\linewidth}p{.15\linewidth}}
    \# & $\;$A & \multicolumn{4}{c}{Game Master} & $\;\:$B\\
    \hline

    \theutterance \stepcounter{utterance}  
    & & & \multicolumn{2}{p{0.3\linewidth}}{
        \cellcolor[rgb]{0.9,0.9,0.9}{
            \makecell[{{p{\linewidth}}}]{
                \texttt{\tiny{[GM$|$GM]}}
                {'hidden': {'team': ['день', 'раковина', 'полиция', 'ягода', 'фильм', 'авария', 'крот', 'яблоко', 'машина'], 'innocent': ['вилка', 'шайба', 'рот', 'лимон', 'амазонка', 'переключатель', 'труба', 'кенгуру'], 'opponent': ['солдат', 'холодный', 'батарея', 'роман', 'грудь', 'африка', 'ночь', 'мексика'], 'assassin': []}, 'revealed': {'team': {'team': [], 'innocent': [], 'opponent': [], 'assassin': []}, 'opponent': {'team': [], 'innocent': [], 'opponent': [], 'assassin': []}}}
            }
        }
    }
    & & \\ \\

    \theutterance \stepcounter{utterance}  
    & & \multicolumn{4}{p{0.6\linewidth}}{
        \cellcolor[rgb]{0.9,0.9,0.9}{
            \makecell[{{p{\linewidth}}}]{
                \texttt{\tiny{[P1$\langle$GM]}}
                \texttt{Дайте одну подсказку, состоящую из одного слова, которая относится к некоторым словам вашей команды из списка ниже. Вы можете выбрать, к скольким и к каким именно словам относится ваша подсказка.} \\
\texttt{Цель — чтобы ваш партнёр угадал все слова вашей команды на основе подсказок, которые вы ему даёте, прежде чем другая команда раскроет все свои. Если ваш партнёр угадает слово 'соперника' или слово{-}'отвлекатор', его ход угадывания заканчивается, и команда соперников получает право сделать ход. Если ваш партнёр угадает слово{-}'убийца', игра заканчивается, и ваша команда немедленно проигрывает.} \\
\texttt{Чем к большему числу слов относится ваша подсказка, тем быстрее вы можете выиграть игру, но тем сложнее вашему партнёру может быть угадать ассоциации, поэтому выбирайте слово подсказки с умом. Слово подсказки должно быть семантически связано с целевыми словами; оно не может быть одним из слов в списках или содержать их части.} \\
\\ 
\texttt{Всегда давайте вашу однословную подсказку и ваш список связанных целевых слов, разделённых запятыми, в следующем формате, и делайте ответы как можно короче; никогда не включайте никакого другого текста, кроме того, что требуется в этой форме:} \\
\\ 
\texttt{ПОДСКАЗКА: <СЛОВО>} \\
\texttt{ЦЕЛИ: <ЦЕЛИ>} \\
\\ 
\texttt{Слова вашей команды: день, раковина, полиция, ягода, фильм, авария, крот, яблоко, машина.} \\
\texttt{Слова команды соперника: солдат, холодный, батарея, роман, грудь, африка, ночь, мексика.} \\
\texttt{Слова{-}отвлекаторы: вилка, шайба, рот, лимон, амазонка, переключатель, труба, кенгуру.} \\
\texttt{Слова{-}убийцы: .} \\
            }
        }
    }
    & \\ \\

    \theutterance \stepcounter{utterance}  
    & \multicolumn{4}{p{0.6\linewidth}}{
        \cellcolor[rgb]{0.9,0.9,0.9}{
            \makecell[{{p{\linewidth}}}]{
                \texttt{\tiny{[P1$\rangle$GM]}}
                \texttt{ПОДСКАЗКА: ДТП} \\
\texttt{ЦЕЛИ: полиция, авария, машина} \\
            }
        }
    }
    & & \\ \\

    \theutterance \stepcounter{utterance}  
    & & & \multicolumn{2}{p{0.3\linewidth}}{
        \cellcolor[rgb]{0.9,0.9,0.9}{
            \makecell[{{p{\linewidth}}}]{
                \texttt{\tiny{[GM$|$GM]}}
                \texttt{дтп} \\
            }
        }
    }
    & & \\ \\

    \theutterance \stepcounter{utterance}  
    & & & \multicolumn{2}{p{0.3\linewidth}}{
        \cellcolor[rgb]{0.9,0.9,0.9}{
            \makecell[{{p{\linewidth}}}]{
                \texttt{\tiny{[GM$|$GM]}}
                ['полиция', 'авария', 'машина']
            }
        }
    }
    & & \\ \\

    \theutterance \stepcounter{utterance}  
    & & & \multicolumn{2}{p{0.3\linewidth}}{
        \cellcolor[rgb]{0.9,0.9,0.9}{
            \makecell[{{p{\linewidth}}}]{
                \texttt{\tiny{[GM$|$GM]}}
                {'word': 'полиция', 'assignment': 'team'}
            }
        }
    }
    & & \\ \\

    \theutterance \stepcounter{utterance}  
    & & & \multicolumn{2}{p{0.3\linewidth}}{
        \cellcolor[rgb]{0.9,0.9,0.9}{
            \makecell[{{p{\linewidth}}}]{
                \texttt{\tiny{[GM$|$GM]}}
                {'word': 'авария', 'assignment': 'team'}
            }
        }
    }
    & & \\ \\

    \theutterance \stepcounter{utterance}  
    & & & \multicolumn{2}{p{0.3\linewidth}}{
        \cellcolor[rgb]{0.9,0.9,0.9}{
            \makecell[{{p{\linewidth}}}]{
                \texttt{\tiny{[GM$|$GM]}}
                {'word': 'машина', 'assignment': 'team'}
            }
        }
    }
    & & \\ \\

    \theutterance \stepcounter{utterance}  
    & & \multicolumn{4}{p{0.6\linewidth}}{
        \cellcolor[rgb]{0.9,0.9,0.9}{
            \makecell[{{p{\linewidth}}}]{
                \texttt{\tiny{[GM$\rangle$P2]}}
                \texttt{Предоставьте список из не более чем 3 слов, разделённых запятыми, из следующего списка, которые лучше всего относятся или наиболее тесно связаны со словом «дтп». Всегда начинайте список вашей(их) догадки(ок) с ОТВЕТ:  и не включайте в свой ответ никакого другого текста.} \\
\\ 
\texttt{авария, холодный, полиция, рот, лимон, шайба, роман, батарея, амазонка, грудь, ночь, день, яблоко, фильм, кенгуру, переключатель, вилка, крот, солдат, мексика, африка, машина, раковина, ягода, труба} \\
            }
        }
    }
    & \\ \\

    \theutterance \stepcounter{utterance}  
    & & & \multicolumn{4}{p{0.6\linewidth}}{
        \cellcolor[rgb]{0.9,0.9,0.9}{
            \makecell[{{p{\linewidth}}}]{
                \texttt{\tiny{[GM$\langle$P2]}}
                \texttt{ОТВЕТ: авария, машина, полиция} \\
            }
        }
    }
     \\ \\

    \theutterance \stepcounter{utterance}  
    & & & \multicolumn{2}{p{0.3\linewidth}}{
        \cellcolor[rgb]{0.9,0.9,0.9}{
            \makecell[{{p{\linewidth}}}]{
                \texttt{\tiny{[GM$|$GM]}}
                ['авария', 'машина', 'полиция']
            }
        }
    }
    & & \\ \\

    \theutterance \stepcounter{utterance}  
    & & & \multicolumn{2}{p{0.3\linewidth}}{
        \cellcolor[rgb]{0.9,0.9,0.9}{
            \makecell[{{p{\linewidth}}}]{
                \texttt{\tiny{[GM$|$GM]}}
                {'word': 'авария', 'assignment': 'team'}
            }
        }
    }
    & & \\ \\

    \theutterance \stepcounter{utterance}  
    & & & \multicolumn{2}{p{0.3\linewidth}}{
        \cellcolor[rgb]{0.9,0.9,0.9}{
            \makecell[{{p{\linewidth}}}]{
                \texttt{\tiny{[GM$|$GM]}}
                {'word': 'авария', 'assignment': 'team'}
            }
        }
    }
    & & \\ \\

    \theutterance \stepcounter{utterance}  
    & & & \multicolumn{2}{p{0.3\linewidth}}{
        \cellcolor[rgb]{0.9,0.9,0.9}{
            \makecell[{{p{\linewidth}}}]{
                \texttt{\tiny{[GM$|$GM]}}
                {'word': 'машина', 'assignment': 'team'}
            }
        }
    }
    & & \\ \\

    \theutterance \stepcounter{utterance}  
    & & & \multicolumn{2}{p{0.3\linewidth}}{
        \cellcolor[rgb]{0.9,0.9,0.9}{
            \makecell[{{p{\linewidth}}}]{
                \texttt{\tiny{[GM$|$GM]}}
                {'word': 'машина', 'assignment': 'team'}
            }
        }
    }
    & & \\ \\

    \theutterance \stepcounter{utterance}  
    & & & \multicolumn{2}{p{0.3\linewidth}}{
        \cellcolor[rgb]{0.9,0.9,0.9}{
            \makecell[{{p{\linewidth}}}]{
                \texttt{\tiny{[GM$|$GM]}}
                {'word': 'полиция', 'assignment': 'team'}
            }
        }
    }
    & & \\ \\

    \theutterance \stepcounter{utterance}  
    & & & \multicolumn{2}{p{0.3\linewidth}}{
        \cellcolor[rgb]{0.9,0.9,0.9}{
            \makecell[{{p{\linewidth}}}]{
                \texttt{\tiny{[GM$|$GM]}}
                {'word': 'полиция', 'assignment': 'team'}
            }
        }
    }
    & & \\ \\

    \theutterance \stepcounter{utterance}  
    & & & \multicolumn{2}{p{0.3\linewidth}}{
        \cellcolor[rgb]{0.9,0.9,0.9}{
            \makecell[{{p{\linewidth}}}]{
                \texttt{\tiny{[GM$|$GM]}}
                {'word': 'солдат', 'assignment': 'opponent'}
            }
        }
    }
    & & \\ \\

    \theutterance \stepcounter{utterance}  
    & & & \multicolumn{2}{p{0.3\linewidth}}{
        \cellcolor[rgb]{0.9,0.9,0.9}{
            \makecell[{{p{\linewidth}}}]{
                \texttt{\tiny{[GM$|$GM]}}
                {'hidden': {'team': ['день', 'раковина', 'ягода', 'фильм', 'крот', 'яблоко'], 'innocent': ['вилка', 'шайба', 'рот', 'лимон', 'амазонка', 'переключатель', 'труба', 'кенгуру'], 'opponent': ['холодный', 'батарея', 'роман', 'грудь', 'африка', 'ночь', 'мексика'], 'assassin': []}, 'revealed': {'team': {'team': ['авария', 'машина', 'полиция'], 'innocent': [], 'opponent': [], 'assassin': []}, 'opponent': {'team': [], 'innocent': [], 'opponent': ['солдат'], 'assassin': []}}}
            }
        }
    }
    & & \\ \\

    \theutterance \stepcounter{utterance}  
    & & \multicolumn{4}{p{0.6\linewidth}}{
        \cellcolor[rgb]{0.9,0.9,0.9}{
            \makecell[{{p{\linewidth}}}]{
                \texttt{\tiny{[P1$\langle$GM]}}
                \texttt{Слова авария, машина, полиция были угаданы правильно. Ход вашего товарища по команде закончился.} \\
\texttt{Теперь предложите ещё одну подсказку, связанную с некоторыми из оставшихся слов вашей команды, и список связанных целевых слов. Помните: начните вашу подсказку с "ПОДСКАЗКА: ", поставьте новую строку и начните ваш список целевых слов, разделённых запятыми, с "ЦЕЛИ: ". Примечание: по сравнению с предыдущими запросами некоторые слова были удалены из списков.} \\
\\ 
\texttt{Оставшиеся слова вашей команды: день, раковина, ягода, фильм, крот, яблоко.} \\
\texttt{Оставшиеся слова вашего соперника: солдат, холодный, батарея, роман, грудь, африка, ночь, мексика.} \\
\texttt{Оставшиеся слова{-}отвлекатели: вилка, шайба, рот, лимон, амазонка, переключатель, труба, кенгуру.} \\
\texttt{Оставшиеся слова{-}убийцы: .} \\
            }
        }
    }
    & \\ \\

    \theutterance \stepcounter{utterance}  
    & \multicolumn{4}{p{0.6\linewidth}}{
        \cellcolor[rgb]{0.9,0.9,0.9}{
            \makecell[{{p{\linewidth}}}]{
                \texttt{\tiny{[P1$\rangle$GM]}}
                \texttt{ПОДСКАЗКА: фрукт} \\
\texttt{ЦЕЛИ: ягода, яблоко} \\
            }
        }
    }
    & & \\ \\

    \theutterance \stepcounter{utterance}  
    & & & \multicolumn{2}{p{0.3\linewidth}}{
        \cellcolor[rgb]{0.9,0.9,0.9}{
            \makecell[{{p{\linewidth}}}]{
                \texttt{\tiny{[GM$|$GM]}}
                \texttt{фрукт} \\
            }
        }
    }
    & & \\ \\

    \theutterance \stepcounter{utterance}  
    & & & \multicolumn{2}{p{0.3\linewidth}}{
        \cellcolor[rgb]{0.9,0.9,0.9}{
            \makecell[{{p{\linewidth}}}]{
                \texttt{\tiny{[GM$|$GM]}}
                ['ягода', 'яблоко']
            }
        }
    }
    & & \\ \\

    \theutterance \stepcounter{utterance}  
    & & & \multicolumn{2}{p{0.3\linewidth}}{
        \cellcolor[rgb]{0.9,0.9,0.9}{
            \makecell[{{p{\linewidth}}}]{
                \texttt{\tiny{[GM$|$GM]}}
                {'word': 'ягода', 'assignment': 'team'}
            }
        }
    }
    & & \\ \\

    \theutterance \stepcounter{utterance}  
    & & & \multicolumn{2}{p{0.3\linewidth}}{
        \cellcolor[rgb]{0.9,0.9,0.9}{
            \makecell[{{p{\linewidth}}}]{
                \texttt{\tiny{[GM$|$GM]}}
                {'word': 'яблоко', 'assignment': 'team'}
            }
        }
    }
    & & \\ \\

    \theutterance \stepcounter{utterance}  
    & & \multicolumn{4}{p{0.6\linewidth}}{
        \cellcolor[rgb]{0.9,0.9,0.9}{
            \makecell[{{p{\linewidth}}}]{
                \texttt{\tiny{[GM$\rangle$P2]}}
                \texttt{Слова авария, машина, полиция были угаданы правильно. Ваш ход закончился.} \\
\texttt{Теперь предоставьте ещё один список, разделённый запятыми, содержащий как минимум 1 и не более чем 2 слов из следующего списка слов, которые лучше всего относятся или наиболее тесно связаны со словом 'фрукт'. Помните, что нужно начать ответ с ОТВЕТ: . Обратите внимание: некоторые слова были удалены из списка по сравнению с предыдущими запросами.} \\
\\ 
\texttt{холодный, рот, лимон, шайба, роман, батарея, амазонка, грудь, ночь, день, яблоко, фильм, кенгуру, переключатель, вилка, крот, мексика, африка, раковина, ягода, труба} \\
            }
        }
    }
    & \\ \\

    \theutterance \stepcounter{utterance}  
    & & & \multicolumn{4}{p{0.6\linewidth}}{
        \cellcolor[rgb]{0.9,0.9,0.9}{
            \makecell[{{p{\linewidth}}}]{
                \texttt{\tiny{[GM$\langle$P2]}}
                \texttt{ОТВЕТ: лимон, яблоко} \\
            }
        }
    }
     \\ \\

    \theutterance \stepcounter{utterance}  
    & & & \multicolumn{2}{p{0.3\linewidth}}{
        \cellcolor[rgb]{0.9,0.9,0.9}{
            \makecell[{{p{\linewidth}}}]{
                \texttt{\tiny{[GM$|$GM]}}
                ['лимон', 'яблоко']
            }
        }
    }
    & & \\ \\

    \theutterance \stepcounter{utterance}  
    & & & \multicolumn{2}{p{0.3\linewidth}}{
        \cellcolor[rgb]{0.9,0.9,0.9}{
            \makecell[{{p{\linewidth}}}]{
                \texttt{\tiny{[GM$|$GM]}}
                {'word': 'лимон', 'assignment': 'innocent'}
            }
        }
    }
    & & \\ \\

    \theutterance \stepcounter{utterance}  
    & & & \multicolumn{2}{p{0.3\linewidth}}{
        \cellcolor[rgb]{0.9,0.9,0.9}{
            \makecell[{{p{\linewidth}}}]{
                \texttt{\tiny{[GM$|$GM]}}
                \texttt{лимон} \\
            }
        }
    }
    & & \\ \\

    \theutterance \stepcounter{utterance}  
    & & & \multicolumn{2}{p{0.3\linewidth}}{
        \cellcolor[rgb]{0.9,0.9,0.9}{
            \makecell[{{p{\linewidth}}}]{
                \texttt{\tiny{[GM$|$GM]}}
                {'word': 'холодный', 'assignment': 'opponent'}
            }
        }
    }
    & & \\ \\

    \theutterance \stepcounter{utterance}  
    & & & \multicolumn{2}{p{0.3\linewidth}}{
        \cellcolor[rgb]{0.9,0.9,0.9}{
            \makecell[{{p{\linewidth}}}]{
                \texttt{\tiny{[GM$|$GM]}}
                {'hidden': {'team': ['день', 'раковина', 'ягода', 'фильм', 'крот', 'яблоко'], 'innocent': ['вилка', 'шайба', 'рот', 'амазонка', 'переключатель', 'труба', 'кенгуру'], 'opponent': ['батарея', 'роман', 'грудь', 'африка', 'ночь', 'мексика'], 'assassin': []}, 'revealed': {'team': {'team': ['авария', 'машина', 'полиция'], 'innocent': ['лимон'], 'opponent': [], 'assassin': []}, 'opponent': {'team': [], 'innocent': [], 'opponent': ['солдат', 'холодный'], 'assassin': []}}}
            }
        }
    }
    & & \\ \\

    \theutterance \stepcounter{utterance}  
    & & \multicolumn{4}{p{0.6\linewidth}}{
        \cellcolor[rgb]{0.9,0.9,0.9}{
            \makecell[{{p{\linewidth}}}]{
                \texttt{\tiny{[P1$\langle$GM]}}
                \texttt{Слово лимон было угадано, но это слово нейтральный. Ход вашего товарища по команде закончился.} \\
\texttt{Теперь предложите ещё одну подсказку, связанную с некоторыми из оставшихся слов вашей команды, и список связанных целевых слов. Помните: начните вашу подсказку с "ПОДСКАЗКА: ", поставьте новую строку и начните ваш список целевых слов, разделённых запятыми, с "ЦЕЛИ: ". Примечание: по сравнению с предыдущими запросами некоторые слова были удалены из списков.} \\
\\ 
\texttt{Оставшиеся слова вашей команды: день, раковина, ягода, фильм, крот, яблоко.} \\
\texttt{Оставшиеся слова вашего соперника: холодный, батарея, роман, грудь, африка, ночь, мексика.} \\
\texttt{Оставшиеся слова{-}отвлекатели: вилка, шайба, рот, амазонка, переключатель, труба, кенгуру.} \\
\texttt{Оставшиеся слова{-}убийцы: .} \\
            }
        }
    }
    & \\ \\

    \theutterance \stepcounter{utterance}  
    & \multicolumn{4}{p{0.6\linewidth}}{
        \cellcolor[rgb]{0.9,0.9,0.9}{
            \makecell[{{p{\linewidth}}}]{
                \texttt{\tiny{[P1$\rangle$GM]}}
                \texttt{ПОДСКАЗКА: садовый} \\
\texttt{ЦЕЛИ: ягода, яблоко} \\
            }
        }
    }
    & & \\ \\

    \theutterance \stepcounter{utterance}  
    & & & \multicolumn{2}{p{0.3\linewidth}}{
        \cellcolor[rgb]{0.9,0.9,0.9}{
            \makecell[{{p{\linewidth}}}]{
                \texttt{\tiny{[GM$|$GM]}}
                \texttt{садовый} \\
            }
        }
    }
    & & \\ \\

    \theutterance \stepcounter{utterance}  
    & & & \multicolumn{2}{p{0.3\linewidth}}{
        \cellcolor[rgb]{0.9,0.9,0.9}{
            \makecell[{{p{\linewidth}}}]{
                \texttt{\tiny{[GM$|$GM]}}
                ['ягода', 'яблоко']
            }
        }
    }
    & & \\ \\

    \theutterance \stepcounter{utterance}  
    & & & \multicolumn{2}{p{0.3\linewidth}}{
        \cellcolor[rgb]{0.9,0.9,0.9}{
            \makecell[{{p{\linewidth}}}]{
                \texttt{\tiny{[GM$|$GM]}}
                {'word': 'ягода', 'assignment': 'team'}
            }
        }
    }
    & & \\ \\

    \theutterance \stepcounter{utterance}  
    & & & \multicolumn{2}{p{0.3\linewidth}}{
        \cellcolor[rgb]{0.9,0.9,0.9}{
            \makecell[{{p{\linewidth}}}]{
                \texttt{\tiny{[GM$|$GM]}}
                {'word': 'яблоко', 'assignment': 'team'}
            }
        }
    }
    & & \\ \\

    \theutterance \stepcounter{utterance}  
    & & \multicolumn{4}{p{0.6\linewidth}}{
        \cellcolor[rgb]{0.9,0.9,0.9}{
            \makecell[{{p{\linewidth}}}]{
                \texttt{\tiny{[GM$\rangle$P2]}}
                \texttt{Слово лимон было угадано, но это слово нейтральный. Ваш ход закончился.} \\
\texttt{Теперь предоставьте ещё один список, разделённый запятыми, содержащий как минимум 1 и не более чем 2 слов из следующего списка слов, которые лучше всего относятся или наиболее тесно связаны со словом 'садовый'. Помните, что нужно начать ответ с ОТВЕТ: . Обратите внимание: некоторые слова были удалены из списка по сравнению с предыдущими запросами.} \\
\\ 
\texttt{рот, шайба, роман, батарея, амазонка, грудь, ночь, день, яблоко, фильм, кенгуру, переключатель, вилка, крот, мексика, африка, раковина, ягода, труба} \\
            }
        }
    }
    & \\ \\

    \theutterance \stepcounter{utterance}  
    & & & \multicolumn{4}{p{0.6\linewidth}}{
        \cellcolor[rgb]{0.9,0.9,0.9}{
            \makecell[{{p{\linewidth}}}]{
                \texttt{\tiny{[GM$\langle$P2]}}
                \texttt{ОТВЕТ: крот, ягода} \\
            }
        }
    }
     \\ \\

    \theutterance \stepcounter{utterance}  
    & & & \multicolumn{2}{p{0.3\linewidth}}{
        \cellcolor[rgb]{0.9,0.9,0.9}{
            \makecell[{{p{\linewidth}}}]{
                \texttt{\tiny{[GM$|$GM]}}
                ['крот', 'ягода']
            }
        }
    }
    & & \\ \\

    \theutterance \stepcounter{utterance}  
    & & & \multicolumn{2}{p{0.3\linewidth}}{
        \cellcolor[rgb]{0.9,0.9,0.9}{
            \makecell[{{p{\linewidth}}}]{
                \texttt{\tiny{[GM$|$GM]}}
                {'word': 'крот', 'assignment': 'team'}
            }
        }
    }
    & & \\ \\

    \theutterance \stepcounter{utterance}  
    & & & \multicolumn{2}{p{0.3\linewidth}}{
        \cellcolor[rgb]{0.9,0.9,0.9}{
            \makecell[{{p{\linewidth}}}]{
                \texttt{\tiny{[GM$|$GM]}}
                {'word': 'ягода', 'assignment': 'team'}
            }
        }
    }
    & & \\ \\

    \theutterance \stepcounter{utterance}  
    & & & \multicolumn{2}{p{0.3\linewidth}}{
        \cellcolor[rgb]{0.9,0.9,0.9}{
            \makecell[{{p{\linewidth}}}]{
                \texttt{\tiny{[GM$|$GM]}}
                {'word': 'ягода', 'assignment': 'team'}
            }
        }
    }
    & & \\ \\

    \theutterance \stepcounter{utterance}  
    & & & \multicolumn{2}{p{0.3\linewidth}}{
        \cellcolor[rgb]{0.9,0.9,0.9}{
            \makecell[{{p{\linewidth}}}]{
                \texttt{\tiny{[GM$|$GM]}}
                {'word': 'роман', 'assignment': 'opponent'}
            }
        }
    }
    & & \\ \\

    \theutterance \stepcounter{utterance}  
    & & & \multicolumn{2}{p{0.3\linewidth}}{
        \cellcolor[rgb]{0.9,0.9,0.9}{
            \makecell[{{p{\linewidth}}}]{
                \texttt{\tiny{[GM$|$GM]}}
                {'hidden': {'team': ['день', 'раковина', 'фильм', 'яблоко'], 'innocent': ['вилка', 'шайба', 'рот', 'амазонка', 'переключатель', 'труба', 'кенгуру'], 'opponent': ['батарея', 'грудь', 'африка', 'ночь', 'мексика'], 'assassin': []}, 'revealed': {'team': {'team': ['авария', 'машина', 'полиция', 'крот', 'ягода'], 'innocent': ['лимон'], 'opponent': [], 'assassin': []}, 'opponent': {'team': [], 'innocent': [], 'opponent': ['солдат', 'холодный', 'роман'], 'assassin': []}}}
            }
        }
    }
    & & \\ \\

    \theutterance \stepcounter{utterance}  
    & & \multicolumn{4}{p{0.6\linewidth}}{
        \cellcolor[rgb]{0.9,0.9,0.9}{
            \makecell[{{p{\linewidth}}}]{
                \texttt{\tiny{[P1$\langle$GM]}}
                \texttt{Слова крот, ягода были угаданы правильно. Ход вашего товарища по команде закончился.} \\
\texttt{Теперь предложите ещё одну подсказку, связанную с некоторыми из оставшихся слов вашей команды, и список связанных целевых слов. Помните: начните вашу подсказку с "ПОДСКАЗКА: ", поставьте новую строку и начните ваш список целевых слов, разделённых запятыми, с "ЦЕЛИ: ". Примечание: по сравнению с предыдущими запросами некоторые слова были удалены из списков.} \\
\\ 
\texttt{Оставшиеся слова вашей команды: день, раковина, фильм, яблоко.} \\
\texttt{Оставшиеся слова вашего соперника: батарея, роман, грудь, африка, ночь, мексика.} \\
\texttt{Оставшиеся слова{-}отвлекатели: вилка, шайба, рот, амазонка, переключатель, труба, кенгуру.} \\
\texttt{Оставшиеся слова{-}убийцы: .} \\
            }
        }
    }
    & \\ \\

    \theutterance \stepcounter{utterance}  
    & \multicolumn{4}{p{0.6\linewidth}}{
        \cellcolor[rgb]{0.9,0.9,0.9}{
            \makecell[{{p{\linewidth}}}]{
                \texttt{\tiny{[P1$\rangle$GM]}}
                \texttt{ПОДСКАЗКА: кино} \\
\texttt{ЦЕЛИ: фильм} \\
            }
        }
    }
    & & \\ \\

    \theutterance \stepcounter{utterance}  
    & & & \multicolumn{2}{p{0.3\linewidth}}{
        \cellcolor[rgb]{0.9,0.9,0.9}{
            \makecell[{{p{\linewidth}}}]{
                \texttt{\tiny{[GM$|$GM]}}
                \texttt{кино} \\
            }
        }
    }
    & & \\ \\

    \theutterance \stepcounter{utterance}  
    & & & \multicolumn{2}{p{0.3\linewidth}}{
        \cellcolor[rgb]{0.9,0.9,0.9}{
            \makecell[{{p{\linewidth}}}]{
                \texttt{\tiny{[GM$|$GM]}}
                ['фильм']
            }
        }
    }
    & & \\ \\

    \theutterance \stepcounter{utterance}  
    & & & \multicolumn{2}{p{0.3\linewidth}}{
        \cellcolor[rgb]{0.9,0.9,0.9}{
            \makecell[{{p{\linewidth}}}]{
                \texttt{\tiny{[GM$|$GM]}}
                {'word': 'фильм', 'assignment': 'team'}
            }
        }
    }
    & & \\ \\

    \theutterance \stepcounter{utterance}  
    & & \multicolumn{4}{p{0.6\linewidth}}{
        \cellcolor[rgb]{0.9,0.9,0.9}{
            \makecell[{{p{\linewidth}}}]{
                \texttt{\tiny{[GM$\rangle$P2]}}
                \texttt{Слова крот, ягода были угаданы правильно. Ваш ход закончился.} \\
\texttt{Теперь предоставьте ещё один список, разделённый запятыми, содержащий как минимум 1 и не более чем 1 слов из следующего списка слов, которые лучше всего относятся или наиболее тесно связаны со словом 'кино'. Помните, что нужно начать ответ с ОТВЕТ: . Обратите внимание: некоторые слова были удалены из списка по сравнению с предыдущими запросами.} \\
\\ 
\texttt{рот, шайба, батарея, амазонка, грудь, ночь, день, яблоко, фильм, кенгуру, переключатель, вилка, мексика, африка, раковина, труба} \\
            }
        }
    }
    & \\ \\

    \theutterance \stepcounter{utterance}  
    & & & \multicolumn{4}{p{0.6\linewidth}}{
        \cellcolor[rgb]{0.9,0.9,0.9}{
            \makecell[{{p{\linewidth}}}]{
                \texttt{\tiny{[GM$\langle$P2]}}
                \texttt{ОТВЕТ: фильм} \\
            }
        }
    }
     \\ \\

    \theutterance \stepcounter{utterance}  
    & & & \multicolumn{2}{p{0.3\linewidth}}{
        \cellcolor[rgb]{0.9,0.9,0.9}{
            \makecell[{{p{\linewidth}}}]{
                \texttt{\tiny{[GM$|$GM]}}
                ['фильм']
            }
        }
    }
    & & \\ \\

    \theutterance \stepcounter{utterance}  
    & & & \multicolumn{2}{p{0.3\linewidth}}{
        \cellcolor[rgb]{0.9,0.9,0.9}{
            \makecell[{{p{\linewidth}}}]{
                \texttt{\tiny{[GM$|$GM]}}
                {'word': 'фильм', 'assignment': 'team'}
            }
        }
    }
    & & \\ \\

    \theutterance \stepcounter{utterance}  
    & & & \multicolumn{2}{p{0.3\linewidth}}{
        \cellcolor[rgb]{0.9,0.9,0.9}{
            \makecell[{{p{\linewidth}}}]{
                \texttt{\tiny{[GM$|$GM]}}
                {'word': 'фильм', 'assignment': 'team'}
            }
        }
    }
    & & \\ \\

    \theutterance \stepcounter{utterance}  
    & & & \multicolumn{2}{p{0.3\linewidth}}{
        \cellcolor[rgb]{0.9,0.9,0.9}{
            \makecell[{{p{\linewidth}}}]{
                \texttt{\tiny{[GM$|$GM]}}
                {'word': 'грудь', 'assignment': 'opponent'}
            }
        }
    }
    & & \\ \\

    \theutterance \stepcounter{utterance}  
    & & & \multicolumn{2}{p{0.3\linewidth}}{
        \cellcolor[rgb]{0.9,0.9,0.9}{
            \makecell[{{p{\linewidth}}}]{
                \texttt{\tiny{[GM$|$GM]}}
                {'hidden': {'team': ['день', 'раковина', 'яблоко'], 'innocent': ['вилка', 'шайба', 'рот', 'амазонка', 'переключатель', 'труба', 'кенгуру'], 'opponent': ['батарея', 'африка', 'ночь', 'мексика'], 'assassin': []}, 'revealed': {'team': {'team': ['авария', 'машина', 'полиция', 'крот', 'ягода', 'фильм'], 'innocent': ['лимон'], 'opponent': [], 'assassin': []}, 'opponent': {'team': [], 'innocent': [], 'opponent': ['солдат', 'холодный', 'роман', 'грудь'], 'assassin': []}}}
            }
        }
    }
    & & \\ \\

    \theutterance \stepcounter{utterance}  
    & & \multicolumn{4}{p{0.6\linewidth}}{
        \cellcolor[rgb]{0.9,0.9,0.9}{
            \makecell[{{p{\linewidth}}}]{
                \texttt{\tiny{[P1$\langle$GM]}}
                \texttt{Слово фильм было угадано правильно. Ход вашего товарища по команде закончился.} \\
\texttt{Теперь предложите ещё одну подсказку, связанную с некоторыми из оставшихся слов вашей команды, и список связанных целевых слов. Помните: начните вашу подсказку с "ПОДСКАЗКА: ", поставьте новую строку и начните ваш список целевых слов, разделённых запятыми, с "ЦЕЛИ: ". Примечание: по сравнению с предыдущими запросами некоторые слова были удалены из списков.} \\
\\ 
\texttt{Оставшиеся слова вашей команды: день, раковина, яблоко.} \\
\texttt{Оставшиеся слова вашего соперника: батарея, грудь, африка, ночь, мексика.} \\
\texttt{Оставшиеся слова{-}отвлекатели: вилка, шайба, рот, амазонка, переключатель, труба, кенгуру.} \\
\texttt{Оставшиеся слова{-}убийцы: .} \\
            }
        }
    }
    & \\ \\

    \theutterance \stepcounter{utterance}  
    & \multicolumn{4}{p{0.6\linewidth}}{
        \cellcolor[rgb]{0.9,0.9,0.9}{
            \makecell[{{p{\linewidth}}}]{
                \texttt{\tiny{[P1$\rangle$GM]}}
                \texttt{ПОДСКАЗКА: кухня} \\
\texttt{ЦЕЛИ: раковина, яблоко} \\
            }
        }
    }
    & & \\ \\

    \theutterance \stepcounter{utterance}  
    & & & \multicolumn{2}{p{0.3\linewidth}}{
        \cellcolor[rgb]{0.9,0.9,0.9}{
            \makecell[{{p{\linewidth}}}]{
                \texttt{\tiny{[GM$|$GM]}}
                \texttt{кухня} \\
            }
        }
    }
    & & \\ \\

    \theutterance \stepcounter{utterance}  
    & & & \multicolumn{2}{p{0.3\linewidth}}{
        \cellcolor[rgb]{0.9,0.9,0.9}{
            \makecell[{{p{\linewidth}}}]{
                \texttt{\tiny{[GM$|$GM]}}
                ['раковина', 'яблоко']
            }
        }
    }
    & & \\ \\

    \theutterance \stepcounter{utterance}  
    & & & \multicolumn{2}{p{0.3\linewidth}}{
        \cellcolor[rgb]{0.9,0.9,0.9}{
            \makecell[{{p{\linewidth}}}]{
                \texttt{\tiny{[GM$|$GM]}}
                {'word': 'раковина', 'assignment': 'team'}
            }
        }
    }
    & & \\ \\

    \theutterance \stepcounter{utterance}  
    & & & \multicolumn{2}{p{0.3\linewidth}}{
        \cellcolor[rgb]{0.9,0.9,0.9}{
            \makecell[{{p{\linewidth}}}]{
                \texttt{\tiny{[GM$|$GM]}}
                {'word': 'яблоко', 'assignment': 'team'}
            }
        }
    }
    & & \\ \\

    \theutterance \stepcounter{utterance}  
    & & \multicolumn{4}{p{0.6\linewidth}}{
        \cellcolor[rgb]{0.9,0.9,0.9}{
            \makecell[{{p{\linewidth}}}]{
                \texttt{\tiny{[GM$\rangle$P2]}}
                \texttt{Слово фильм было угадано правильно. Ваш ход закончился.} \\
\texttt{Теперь предоставьте ещё один список, разделённый запятыми, содержащий как минимум 1 и не более чем 2 слов из следующего списка слов, которые лучше всего относятся или наиболее тесно связаны со словом 'кухня'. Помните, что нужно начать ответ с ОТВЕТ: . Обратите внимание: некоторые слова были удалены из списка по сравнению с предыдущими запросами.} \\
\\ 
\texttt{рот, шайба, батарея, амазонка, ночь, день, яблоко, кенгуру, переключатель, вилка, мексика, африка, раковина, труба} \\
            }
        }
    }
    & \\ \\

    \theutterance \stepcounter{utterance}  
    & & & \multicolumn{4}{p{0.6\linewidth}}{
        \cellcolor[rgb]{0.9,0.9,0.9}{
            \makecell[{{p{\linewidth}}}]{
                \texttt{\tiny{[GM$\langle$P2]}}
                \texttt{ОТВЕТ: вилка, раковина} \\
            }
        }
    }
     \\ \\

    \theutterance \stepcounter{utterance}  
    & & & \multicolumn{2}{p{0.3\linewidth}}{
        \cellcolor[rgb]{0.9,0.9,0.9}{
            \makecell[{{p{\linewidth}}}]{
                \texttt{\tiny{[GM$|$GM]}}
                ['вилка', 'раковина']
            }
        }
    }
    & & \\ \\

    \theutterance \stepcounter{utterance}  
    & & & \multicolumn{2}{p{0.3\linewidth}}{
        \cellcolor[rgb]{0.9,0.9,0.9}{
            \makecell[{{p{\linewidth}}}]{
                \texttt{\tiny{[GM$|$GM]}}
                {'word': 'вилка', 'assignment': 'innocent'}
            }
        }
    }
    & & \\ \\

    \theutterance \stepcounter{utterance}  
    & & & \multicolumn{2}{p{0.3\linewidth}}{
        \cellcolor[rgb]{0.9,0.9,0.9}{
            \makecell[{{p{\linewidth}}}]{
                \texttt{\tiny{[GM$|$GM]}}
                \texttt{вилка} \\
            }
        }
    }
    & & \\ \\

    \theutterance \stepcounter{utterance}  
    & & & \multicolumn{2}{p{0.3\linewidth}}{
        \cellcolor[rgb]{0.9,0.9,0.9}{
            \makecell[{{p{\linewidth}}}]{
                \texttt{\tiny{[GM$|$GM]}}
                {'word': 'батарея', 'assignment': 'opponent'}
            }
        }
    }
    & & \\ \\

    \theutterance \stepcounter{utterance}  
    & & & \multicolumn{2}{p{0.3\linewidth}}{
        \cellcolor[rgb]{0.9,0.9,0.9}{
            \makecell[{{p{\linewidth}}}]{
                \texttt{\tiny{[GM$|$GM]}}
                {'hidden': {'team': ['день', 'раковина', 'яблоко'], 'innocent': ['шайба', 'рот', 'амазонка', 'переключатель', 'труба', 'кенгуру'], 'opponent': ['африка', 'ночь', 'мексика'], 'assassin': []}, 'revealed': {'team': {'team': ['авария', 'машина', 'полиция', 'крот', 'ягода', 'фильм'], 'innocent': ['лимон', 'вилка'], 'opponent': [], 'assassin': []}, 'opponent': {'team': [], 'innocent': [], 'opponent': ['солдат', 'холодный', 'роман', 'грудь', 'батарея'], 'assassin': []}}}
            }
        }
    }
    & & \\ \\

    \theutterance \stepcounter{utterance}  
    & & \multicolumn{4}{p{0.6\linewidth}}{
        \cellcolor[rgb]{0.9,0.9,0.9}{
            \makecell[{{p{\linewidth}}}]{
                \texttt{\tiny{[P1$\langle$GM]}}
                \texttt{Слово вилка было угадано, но это слово нейтральный. Ход вашего товарища по команде закончился.} \\
\texttt{Теперь предложите ещё одну подсказку, связанную с некоторыми из оставшихся слов вашей команды, и список связанных целевых слов. Помните: начните вашу подсказку с "ПОДСКАЗКА: ", поставьте новую строку и начните ваш список целевых слов, разделённых запятыми, с "ЦЕЛИ: ". Примечание: по сравнению с предыдущими запросами некоторые слова были удалены из списков.} \\
\\ 
\texttt{Оставшиеся слова вашей команды: день, раковина, яблоко.} \\
\texttt{Оставшиеся слова вашего соперника: батарея, африка, ночь, мексика.} \\
\texttt{Оставшиеся слова{-}отвлекатели: шайба, рот, амазонка, переключатель, труба, кенгуру.} \\
\texttt{Оставшиеся слова{-}убийцы: .} \\
            }
        }
    }
    & \\ \\

    \theutterance \stepcounter{utterance}  
    & \multicolumn{4}{p{0.6\linewidth}}{
        \cellcolor[rgb]{0.9,0.9,0.9}{
            \makecell[{{p{\linewidth}}}]{
                \texttt{\tiny{[P1$\rangle$GM]}}
                \texttt{ПОДСКАЗКА: сантехника} \\
\texttt{ЦЕЛИ: раковина} \\
            }
        }
    }
    & & \\ \\

    \theutterance \stepcounter{utterance}  
    & & & \multicolumn{2}{p{0.3\linewidth}}{
        \cellcolor[rgb]{0.9,0.9,0.9}{
            \makecell[{{p{\linewidth}}}]{
                \texttt{\tiny{[GM$|$GM]}}
                \texttt{сантехника} \\
            }
        }
    }
    & & \\ \\

    \theutterance \stepcounter{utterance}  
    & & & \multicolumn{2}{p{0.3\linewidth}}{
        \cellcolor[rgb]{0.9,0.9,0.9}{
            \makecell[{{p{\linewidth}}}]{
                \texttt{\tiny{[GM$|$GM]}}
                ['раковина']
            }
        }
    }
    & & \\ \\

    \theutterance \stepcounter{utterance}  
    & & & \multicolumn{2}{p{0.3\linewidth}}{
        \cellcolor[rgb]{0.9,0.9,0.9}{
            \makecell[{{p{\linewidth}}}]{
                \texttt{\tiny{[GM$|$GM]}}
                {'word': 'раковина', 'assignment': 'team'}
            }
        }
    }
    & & \\ \\

    \theutterance \stepcounter{utterance}  
    & & \multicolumn{4}{p{0.6\linewidth}}{
        \cellcolor[rgb]{0.9,0.9,0.9}{
            \makecell[{{p{\linewidth}}}]{
                \texttt{\tiny{[GM$\rangle$P2]}}
                \texttt{Слово вилка было угадано, но это слово нейтральный. Ваш ход закончился.} \\
\texttt{Теперь предоставьте ещё один список, разделённый запятыми, содержащий как минимум 1 и не более чем 1 слов из следующего списка слов, которые лучше всего относятся или наиболее тесно связаны со словом 'сантехника'. Помните, что нужно начать ответ с ОТВЕТ: . Обратите внимание: некоторые слова были удалены из списка по сравнению с предыдущими запросами.} \\
\\ 
\texttt{рот, шайба, амазонка, ночь, день, яблоко, кенгуру, переключатель, мексика, африка, раковина, труба} \\
            }
        }
    }
    & \\ \\

    \theutterance \stepcounter{utterance}  
    & & & \multicolumn{4}{p{0.6\linewidth}}{
        \cellcolor[rgb]{0.9,0.9,0.9}{
            \makecell[{{p{\linewidth}}}]{
                \texttt{\tiny{[GM$\langle$P2]}}
                \texttt{ОТВЕТ: труба} \\
            }
        }
    }
     \\ \\

    \theutterance \stepcounter{utterance}  
    & & & \multicolumn{2}{p{0.3\linewidth}}{
        \cellcolor[rgb]{0.9,0.9,0.9}{
            \makecell[{{p{\linewidth}}}]{
                \texttt{\tiny{[GM$|$GM]}}
                ['труба']
            }
        }
    }
    & & \\ \\

    \theutterance \stepcounter{utterance}  
    & & & \multicolumn{2}{p{0.3\linewidth}}{
        \cellcolor[rgb]{0.9,0.9,0.9}{
            \makecell[{{p{\linewidth}}}]{
                \texttt{\tiny{[GM$|$GM]}}
                {'word': 'труба', 'assignment': 'innocent'}
            }
        }
    }
    & & \\ \\

    \theutterance \stepcounter{utterance}  
    & & & \multicolumn{2}{p{0.3\linewidth}}{
        \cellcolor[rgb]{0.9,0.9,0.9}{
            \makecell[{{p{\linewidth}}}]{
                \texttt{\tiny{[GM$|$GM]}}
                \texttt{труба} \\
            }
        }
    }
    & & \\ \\

    \theutterance \stepcounter{utterance}  
    & & & \multicolumn{2}{p{0.3\linewidth}}{
        \cellcolor[rgb]{0.9,0.9,0.9}{
            \makecell[{{p{\linewidth}}}]{
                \texttt{\tiny{[GM$|$GM]}}
                {'word': 'мексика', 'assignment': 'opponent'}
            }
        }
    }
    & & \\ \\

    \theutterance \stepcounter{utterance}  
    & & & \multicolumn{2}{p{0.3\linewidth}}{
        \cellcolor[rgb]{0.9,0.9,0.9}{
            \makecell[{{p{\linewidth}}}]{
                \texttt{\tiny{[GM$|$GM]}}
                {'hidden': {'team': ['день', 'раковина', 'яблоко'], 'innocent': ['шайба', 'рот', 'амазонка', 'переключатель', 'кенгуру'], 'opponent': ['африка', 'ночь'], 'assassin': []}, 'revealed': {'team': {'team': ['авария', 'машина', 'полиция', 'крот', 'ягода', 'фильм'], 'innocent': ['лимон', 'вилка', 'труба'], 'opponent': [], 'assassin': []}, 'opponent': {'team': [], 'innocent': [], 'opponent': ['солдат', 'холодный', 'роман', 'грудь', 'батарея', 'мексика'], 'assassin': []}}}
            }
        }
    }
    & & \\ \\

    \theutterance \stepcounter{utterance}  
    & & \multicolumn{4}{p{0.6\linewidth}}{
        \cellcolor[rgb]{0.9,0.9,0.9}{
            \makecell[{{p{\linewidth}}}]{
                \texttt{\tiny{[P1$\langle$GM]}}
                \texttt{Слово труба было угадано, но это слово нейтральный. Ход вашего товарища по команде закончился.} \\
\texttt{Теперь предложите ещё одну подсказку, связанную с некоторыми из оставшихся слов вашей команды, и список связанных целевых слов. Помните: начните вашу подсказку с "ПОДСКАЗКА: ", поставьте новую строку и начните ваш список целевых слов, разделённых запятыми, с "ЦЕЛИ: ". Примечание: по сравнению с предыдущими запросами некоторые слова были удалены из списков.} \\
\\ 
\texttt{Оставшиеся слова вашей команды: день, раковина, яблоко.} \\
\texttt{Оставшиеся слова вашего соперника: африка, ночь, мексика.} \\
\texttt{Оставшиеся слова{-}отвлекатели: шайба, рот, амазонка, переключатель, кенгуру.} \\
\texttt{Оставшиеся слова{-}убийцы: .} \\
            }
        }
    }
    & \\ \\

    \theutterance \stepcounter{utterance}  
    & \multicolumn{4}{p{0.6\linewidth}}{
        \cellcolor[rgb]{0.9,0.9,0.9}{
            \makecell[{{p{\linewidth}}}]{
                \texttt{\tiny{[P1$\rangle$GM]}}
                \texttt{ПОДСКАЗКА: мойка} \\
\texttt{ЦЕЛИ: раковина} \\
            }
        }
    }
    & & \\ \\

    \theutterance \stepcounter{utterance}  
    & & & \multicolumn{2}{p{0.3\linewidth}}{
        \cellcolor[rgb]{0.9,0.9,0.9}{
            \makecell[{{p{\linewidth}}}]{
                \texttt{\tiny{[GM$|$GM]}}
                \texttt{мойка} \\
            }
        }
    }
    & & \\ \\

    \theutterance \stepcounter{utterance}  
    & & & \multicolumn{2}{p{0.3\linewidth}}{
        \cellcolor[rgb]{0.9,0.9,0.9}{
            \makecell[{{p{\linewidth}}}]{
                \texttt{\tiny{[GM$|$GM]}}
                ['раковина']
            }
        }
    }
    & & \\ \\

    \theutterance \stepcounter{utterance}  
    & & & \multicolumn{2}{p{0.3\linewidth}}{
        \cellcolor[rgb]{0.9,0.9,0.9}{
            \makecell[{{p{\linewidth}}}]{
                \texttt{\tiny{[GM$|$GM]}}
                {'word': 'раковина', 'assignment': 'team'}
            }
        }
    }
    & & \\ \\

    \theutterance \stepcounter{utterance}  
    & & \multicolumn{4}{p{0.6\linewidth}}{
        \cellcolor[rgb]{0.9,0.9,0.9}{
            \makecell[{{p{\linewidth}}}]{
                \texttt{\tiny{[GM$\rangle$P2]}}
                \texttt{Слово труба было угадано, но это слово нейтральный. Ваш ход закончился.} \\
\texttt{Теперь предоставьте ещё один список, разделённый запятыми, содержащий как минимум 1 и не более чем 1 слов из следующего списка слов, которые лучше всего относятся или наиболее тесно связаны со словом 'мойка'. Помните, что нужно начать ответ с ОТВЕТ: . Обратите внимание: некоторые слова были удалены из списка по сравнению с предыдущими запросами.} \\
\\ 
\texttt{рот, шайба, амазонка, ночь, день, яблоко, кенгуру, переключатель, африка, раковина} \\
            }
        }
    }
    & \\ \\

    \theutterance \stepcounter{utterance}  
    & & & \multicolumn{4}{p{0.6\linewidth}}{
        \cellcolor[rgb]{0.9,0.9,0.9}{
            \makecell[{{p{\linewidth}}}]{
                \texttt{\tiny{[GM$\langle$P2]}}
                \texttt{ОТВЕТ: раковина} \\
            }
        }
    }
     \\ \\

    \theutterance \stepcounter{utterance}  
    & & & \multicolumn{2}{p{0.3\linewidth}}{
        \cellcolor[rgb]{0.9,0.9,0.9}{
            \makecell[{{p{\linewidth}}}]{
                \texttt{\tiny{[GM$|$GM]}}
                ['раковина']
            }
        }
    }
    & & \\ \\

    \theutterance \stepcounter{utterance}  
    & & & \multicolumn{2}{p{0.3\linewidth}}{
        \cellcolor[rgb]{0.9,0.9,0.9}{
            \makecell[{{p{\linewidth}}}]{
                \texttt{\tiny{[GM$|$GM]}}
                {'word': 'раковина', 'assignment': 'team'}
            }
        }
    }
    & & \\ \\

    \theutterance \stepcounter{utterance}  
    & & & \multicolumn{2}{p{0.3\linewidth}}{
        \cellcolor[rgb]{0.9,0.9,0.9}{
            \makecell[{{p{\linewidth}}}]{
                \texttt{\tiny{[GM$|$GM]}}
                {'word': 'раковина', 'assignment': 'team'}
            }
        }
    }
    & & \\ \\

    \theutterance \stepcounter{utterance}  
    & & & \multicolumn{2}{p{0.3\linewidth}}{
        \cellcolor[rgb]{0.9,0.9,0.9}{
            \makecell[{{p{\linewidth}}}]{
                \texttt{\tiny{[GM$|$GM]}}
                {'word': 'африка', 'assignment': 'opponent'}
            }
        }
    }
    & & \\ \\

    \theutterance \stepcounter{utterance}  
    & & & \multicolumn{2}{p{0.3\linewidth}}{
        \cellcolor[rgb]{0.9,0.9,0.9}{
            \makecell[{{p{\linewidth}}}]{
                \texttt{\tiny{[GM$|$GM]}}
                {'hidden': {'team': ['день', 'яблоко'], 'innocent': ['шайба', 'рот', 'амазонка', 'переключатель', 'кенгуру'], 'opponent': ['ночь'], 'assassin': []}, 'revealed': {'team': {'team': ['авария', 'машина', 'полиция', 'крот', 'ягода', 'фильм', 'раковина'], 'innocent': ['лимон', 'вилка', 'труба'], 'opponent': [], 'assassin': []}, 'opponent': {'team': [], 'innocent': [], 'opponent': ['солдат', 'холодный', 'роман', 'грудь', 'батарея', 'мексика', 'африка'], 'assassin': []}}}
            }
        }
    }
    & & \\ \\

    \theutterance \stepcounter{utterance}  
    & & \multicolumn{4}{p{0.6\linewidth}}{
        \cellcolor[rgb]{0.9,0.9,0.9}{
            \makecell[{{p{\linewidth}}}]{
                \texttt{\tiny{[P1$\langle$GM]}}
                \texttt{Слово раковина было угадано правильно. Ход вашего товарища по команде закончился.} \\
\texttt{Теперь предложите ещё одну подсказку, связанную с некоторыми из оставшихся слов вашей команды, и список связанных целевых слов. Помните: начните вашу подсказку с "ПОДСКАЗКА: ", поставьте новую строку и начните ваш список целевых слов, разделённых запятыми, с "ЦЕЛИ: ". Примечание: по сравнению с предыдущими запросами некоторые слова были удалены из списков.} \\
\\ 
\texttt{Оставшиеся слова вашей команды: день, яблоко.} \\
\texttt{Оставшиеся слова вашего соперника: африка, ночь.} \\
\texttt{Оставшиеся слова{-}отвлекатели: шайба, рот, амазонка, переключатель, кенгуру.} \\
\texttt{Оставшиеся слова{-}убийцы: .} \\
            }
        }
    }
    & \\ \\

    \theutterance \stepcounter{utterance}  
    & \multicolumn{4}{p{0.6\linewidth}}{
        \cellcolor[rgb]{0.9,0.9,0.9}{
            \makecell[{{p{\linewidth}}}]{
                \texttt{\tiny{[P1$\rangle$GM]}}
                \texttt{ПОДСКАЗКА: сутки} \\
\texttt{ЦЕЛИ: день} \\
            }
        }
    }
    & & \\ \\

    \theutterance \stepcounter{utterance}  
    & & & \multicolumn{2}{p{0.3\linewidth}}{
        \cellcolor[rgb]{0.9,0.9,0.9}{
            \makecell[{{p{\linewidth}}}]{
                \texttt{\tiny{[GM$|$GM]}}
                \texttt{сутки} \\
            }
        }
    }
    & & \\ \\

    \theutterance \stepcounter{utterance}  
    & & & \multicolumn{2}{p{0.3\linewidth}}{
        \cellcolor[rgb]{0.9,0.9,0.9}{
            \makecell[{{p{\linewidth}}}]{
                \texttt{\tiny{[GM$|$GM]}}
                ['день']
            }
        }
    }
    & & \\ \\

    \theutterance \stepcounter{utterance}  
    & & & \multicolumn{2}{p{0.3\linewidth}}{
        \cellcolor[rgb]{0.9,0.9,0.9}{
            \makecell[{{p{\linewidth}}}]{
                \texttt{\tiny{[GM$|$GM]}}
                {'word': 'день', 'assignment': 'team'}
            }
        }
    }
    & & \\ \\

    \theutterance \stepcounter{utterance}  
    & & \multicolumn{4}{p{0.6\linewidth}}{
        \cellcolor[rgb]{0.9,0.9,0.9}{
            \makecell[{{p{\linewidth}}}]{
                \texttt{\tiny{[GM$\rangle$P2]}}
                \texttt{Слово раковина было угадано правильно. Ваш ход закончился.} \\
\texttt{Теперь предоставьте ещё один список, разделённый запятыми, содержащий как минимум 1 и не более чем 1 слов из следующего списка слов, которые лучше всего относятся или наиболее тесно связаны со словом 'сутки'. Помните, что нужно начать ответ с ОТВЕТ: . Обратите внимание: некоторые слова были удалены из списка по сравнению с предыдущими запросами.} \\
\\ 
\texttt{рот, шайба, амазонка, ночь, день, яблоко, кенгуру, переключатель} \\
            }
        }
    }
    & \\ \\

    \theutterance \stepcounter{utterance}  
    & & & \multicolumn{4}{p{0.6\linewidth}}{
        \cellcolor[rgb]{0.9,0.9,0.9}{
            \makecell[{{p{\linewidth}}}]{
                \texttt{\tiny{[GM$\langle$P2]}}
                \texttt{ОТВЕТ: день} \\
            }
        }
    }
     \\ \\

    \theutterance \stepcounter{utterance}  
    & & & \multicolumn{2}{p{0.3\linewidth}}{
        \cellcolor[rgb]{0.9,0.9,0.9}{
            \makecell[{{p{\linewidth}}}]{
                \texttt{\tiny{[GM$|$GM]}}
                ['день']
            }
        }
    }
    & & \\ \\

    \theutterance \stepcounter{utterance}  
    & & & \multicolumn{2}{p{0.3\linewidth}}{
        \cellcolor[rgb]{0.9,0.9,0.9}{
            \makecell[{{p{\linewidth}}}]{
                \texttt{\tiny{[GM$|$GM]}}
                {'word': 'день', 'assignment': 'team'}
            }
        }
    }
    & & \\ \\

    \theutterance \stepcounter{utterance}  
    & & & \multicolumn{2}{p{0.3\linewidth}}{
        \cellcolor[rgb]{0.9,0.9,0.9}{
            \makecell[{{p{\linewidth}}}]{
                \texttt{\tiny{[GM$|$GM]}}
                {'word': 'день', 'assignment': 'team'}
            }
        }
    }
    & & \\ \\

    \theutterance \stepcounter{utterance}  
    & & & \multicolumn{2}{p{0.3\linewidth}}{
        \cellcolor[rgb]{0.9,0.9,0.9}{
            \makecell[{{p{\linewidth}}}]{
                \texttt{\tiny{[GM$|$GM]}}
                {'word': 'ночь', 'assignment': 'opponent'}
            }
        }
    }
    & & \\ \\

    \theutterance \stepcounter{utterance}  
    & & & \multicolumn{2}{p{0.3\linewidth}}{
        \cellcolor[rgb]{0.9,0.9,0.9}{
            \makecell[{{p{\linewidth}}}]{
                \texttt{\tiny{[GM$|$GM]}}
                {'hidden': {'team': ['яблоко'], 'innocent': ['шайба', 'рот', 'амазонка', 'переключатель', 'кенгуру'], 'opponent': [], 'assassin': []}, 'revealed': {'team': {'team': ['авария', 'машина', 'полиция', 'крот', 'ягода', 'фильм', 'раковина', 'день'], 'innocent': ['лимон', 'вилка', 'труба'], 'opponent': [], 'assassin': []}, 'opponent': {'team': [], 'innocent': [], 'opponent': ['солдат', 'холодный', 'роман', 'грудь', 'батарея', 'мексика', 'африка', 'ночь'], 'assassin': []}}}
            }
        }
    }
    & & \\ \\

    \theutterance \stepcounter{utterance}  
    & & \multicolumn{4}{p{0.6\linewidth}}{
        \cellcolor[rgb]{0.9,0.9,0.9}{
            \makecell[{{p{\linewidth}}}]{
                \texttt{\tiny{[P1$\langle$GM]}}
                \texttt{Слово день было угадано правильно. Ход вашего товарища по команде закончился.} \\
\texttt{Теперь предложите ещё одну подсказку, связанную с некоторыми из оставшихся слов вашей команды, и список связанных целевых слов. Помните: начните вашу подсказку с "ПОДСКАЗКА: ", поставьте новую строку и начните ваш список целевых слов, разделённых запятыми, с "ЦЕЛИ: ". Примечание: по сравнению с предыдущими запросами некоторые слова были удалены из списков.} \\
\\ 
\texttt{Оставшиеся слова вашей команды: яблоко.} \\
\texttt{Оставшиеся слова вашего соперника: ночь.} \\
\texttt{Оставшиеся слова{-}отвлекатели: шайба, рот, амазонка, переключатель, кенгуру.} \\
\texttt{Оставшиеся слова{-}убийцы: .} \\
            }
        }
    }
    & \\ \\

    \theutterance \stepcounter{utterance}  
    & \multicolumn{4}{p{0.6\linewidth}}{
        \cellcolor[rgb]{0.9,0.9,0.9}{
            \makecell[{{p{\linewidth}}}]{
                \texttt{\tiny{[P1$\rangle$GM]}}
                \texttt{ПОДСКАЗКА: фрукт} \\
\texttt{ЦЕЛИ: яблоко} \\
            }
        }
    }
    & & \\ \\

    \theutterance \stepcounter{utterance}  
    & & & \multicolumn{2}{p{0.3\linewidth}}{
        \cellcolor[rgb]{0.9,0.9,0.9}{
            \makecell[{{p{\linewidth}}}]{
                \texttt{\tiny{[GM$|$GM]}}
                \texttt{фрукт} \\
            }
        }
    }
    & & \\ \\

    \theutterance \stepcounter{utterance}  
    & & & \multicolumn{2}{p{0.3\linewidth}}{
        \cellcolor[rgb]{0.9,0.9,0.9}{
            \makecell[{{p{\linewidth}}}]{
                \texttt{\tiny{[GM$|$GM]}}
                ['яблоко']
            }
        }
    }
    & & \\ \\

    \theutterance \stepcounter{utterance}  
    & & & \multicolumn{2}{p{0.3\linewidth}}{
        \cellcolor[rgb]{0.9,0.9,0.9}{
            \makecell[{{p{\linewidth}}}]{
                \texttt{\tiny{[GM$|$GM]}}
                {'word': 'яблоко', 'assignment': 'team'}
            }
        }
    }
    & & \\ \\

    \theutterance \stepcounter{utterance}  
    & & & \multicolumn{2}{p{0.3\linewidth}}{
        \cellcolor[rgb]{0.9,0.9,0.9}{
            \makecell[{{p{\linewidth}}}]{
                \texttt{\tiny{[GM$|$GM]}}
                \texttt{opponent has won} \\
            }
        }
    }
    & & \\ \\

\end{supertabular}
}

\end{document}
